\addcontentsline{toc}{chapter}{Abstract}\vspace{-1cm}
%Border
\begin{tikzpicture}[remember picture, overlay]
  \draw[line width = 4pt] ($(current page.north west) + (0.75in,-0.75in)$) rectangle ($(current page.south east) + (-0.75in,0.75in)$);
\end{tikzpicture}

Protecting coastlines and managing vessel traffic are fundamental to modern maritime security. The global market for these surveillance technologies is expanding rapidly, valued at roughly 25 billion USD today and projected to cross 35 billion USD by the early 2030s. However, traditional monitoring tools like radar and manual watch-standing are struggling to keep up. They lack scalability, fail to deliver real-time analytics, and suffer from massive blind spots during heavy fog or sun glare. More importantly, they cannot predict where a ship is heading. To fix this, the maritime industry desperately needs automated, vision-driven solutions.

To address these vulnerabilities, a complete artificial intelligence pipeline was developed specifically for modern maritime surveillance. The system tackles three major operational hurdles: detecting multiple ships simultaneously, maintaining their distinct identities, and forecasting their future paths to prevent collisions. At the core of the framework is a \gls{yolov8} network, which handles real-time vessel detection. Once a boat is spotted, the \gls{deepocsort} algorithm takes over. By pairing a Kalman filter with unique visual appearance embeddings (\gls{reid}), the tracker continuously follows a vessel even through temporary occlusions. Finally, a Long Short-Term Memory (\gls{lstm}) neural network forecasts future trajectories while analytically calculating critical metrics like speed and heading.

The entire architecture was constructed in Python, leveraging \gls{pytorch}, \gls{opencv}, and dedicated \gls{cuda} hardware to sustain real-time processing speeds. When pushed through challenging video tests packed with dense traffic and sudden lighting drops, the system consistently maintained strong detection rates and highly accurate trajectory forecasts. Validated directly on physical \glspl{gpu} to prove its real-time viability, this modular design provides a highly scalable foundation. It stands ready to ingest multi-camera feeds and link with existing \gls{ais} networks, bringing next-generation domain awareness to a production scale.

\pagebreak